\section{Referencia rápida de términos}

\subsection{Tabla de términos clave}

\begin{tabular}{p{0.25\textwidth} p{0.7\textwidth}}
\toprule
Término & Explicación \\
\midrule
Agent Builder & Herramienta visual de Vertex AI para definir el planner, las tools y el feedback loop \\
Policy-as-code & Define las reglas de política mediante código ejecutable, combinándolo con OPA y otras herramientas para el gatekeeping \\
Tool manifest & Describe la entrada, la salida, la latency y el failure mode de cada herramienta \\
Feedback loop & Envía el user feedback y los failure tags al labeling \& retrain pipeline \\
Observability pipeline & Agrega latency, tool usage y alerts en el dashboard \\
\bottomrule
\end{tabular}

\begin{importantbox}{La función del glosario}
Internamente: ayuda a que los ingenieros en onboarding entiendan la plataforma rápidamente; de cara al exterior: le da al equipo de compliance definiciones unificadas. Se recomienda colocar el glosario en el README y darle mantenimiento periódico.
\end{importantbox}

\subsection{Aplicación y recomendaciones de mantenimiento}

\begin{itemize}
    \item Llevar el glossary bajo control de versiones y sincronizarlo con el README, para asegurar que los términos se actualicen en cuanto cambie la policy
    \item Durante el onboarding de los nuevos integrantes, pedirles que repitan el glossary para asegurar la comprensión y la memorización
    \item Convertir el glossary en un `cheatsheet` que se use en las compliance meeting y en el board review
\end{itemize}

